% Pursuing Stacks (cotes 134-2 to 134-8): archive cote 134-5, site page 55 (PDF page 56), typescript page 242.
% typescript. Draft compared: ps4.tex:2771-2844.
% AI-transcribed by Claude Opus 5.5 (claude-opus-5-5, adaptive thinking, effort high), Message Batches API, batch ps2 (msgbatch_01NVwr1Yuo8qJf5rykeGBXPK), 2026-11-12; reviewed against both supplied scan halves; see CHECK_DECISIONS.tsv.
\scanpage{134-5/55}
\mspage{242}
tisfactory for f\unsure{or}mulating the following statement,  which comes out here
rather tautologically,  and which however appears to me as exactly what
I had been looking for in the ``observation'' recalled above:

\hand{\uline{\textcolor{black}{Theorem}1}} Let $(M,M_{as})$  be a \hand{\uline{\textcolor{black}{modelizing}}} asphericity structure,  \xout{and} A a
\xout{smallxtest} pseudo-test category,  and
\[ i:\ A \hand{\longrightarrow} M \]
a functor,  factoring through $M_{as}$  (i.e. $i(a)$ is aspheric for any a in A)\hand{.}
\inserted[typed]{We assume moreover M has a final object $e_M$ .}
Then the six conditions (i) to (ii") of prop.4  (p.236)\hand{,} \hand{(}the first
of which expresses that i is $M_{as}$-aspheric\hand{)} are equivalent,  and they are
equival\unsure{e|a}nt to the following condition:
\xout{(iii)}  (iii)  The functor  $i^{*}$: $M \hand{\to} A\hat{}$ gives rise to a functor
\[ Hot_M = W^{-1}M \hand{\longrightarrow} Hot_A = \underline{W}_A{}^{-1}A\hat{} \]
(i.e.  (ii') of prop.4 holds,  namely  $i^{*}(W) \hand{\subset} \underline{W}_A$ ),  \hand{\uline{\textcolor{black}{and}}}  this functor
moreover is an equivalence of categories.

Proof.  For the first statement (equivalence of conditiohs (i) to
(ii')), \inserted[typed]{By prop.4} we need only show that the asphericity structure  and A are as-%
pheric.  But \xout{itxis} this follows from the assumptions and \inserted{from} the

\hand{\uline{\textcolor{black}{Lemma}}} : Let $(M,M_{as})$  be any modelizing asphericity structure,  then a
final object of M is aspheric\inserted[typed]{(i.e. the structure is ``aspheric'').} In particular,  if A is a pseudo-test
category (i.e.  $(A\hat{},A\hat{}_{as})$  is a modelizing asphericity structure),  then A
is aspheric.

We only have to prove the first statement.  \xout{Let} It is immediate
that $e_M$  is a final object of $W^{-1}M = Hot_M$  (this is valid whenever we
got a localization $W^{-1}M$  of a category M with final object $e_M$ ),  hence
it's image in \xout{(Cat}(Hot) = $\underline{W}^{-1}$(Cat)  is a final object.  Thus we are redu-%
ced to proving the following

\hand{\uline{\textcolor{black}{Corollary}}} :  Let $(M,M_{as})$ be any asphericity structure,  consider the
composition
\[ \hand{\phi}_M\ :\ M \hand{\longrightarrow} Hot_M = W^{-1}M \hand{\xrightarrow{can.}} (Hot) \quad , \]
then we get
\[ M_{as} = \hand{\{} x \hand{\in} Ob\ M \ \hand{|}\ \hand{\phi}_M(x) \text{ is a final object of (Hot)}\hand{\}}. \]

Indeed,  using the construction of  $Hot_M \hand{\longrightarrow} (Hot)$  in terms of a
given $M_{as}$-aspheric functor  $B \hand{\to} M$ ,  we are reduced to the same state-%
ment,  with $(M,M_{as})$  replaced by $(B\hat{},\ B\hat{}_{as})$.  The statement then \over{d}{r}educes
to :  the object B in (Cat)  is aspheric iff it's image in (Hot) is a
final object  (this can be viewed also as the particular case of the
\margin{55}
% Small displaced marks occur below the n in Indeed and above/below the h area in The statement; their precise correction function is unresolved.
